\section{Einleitung}

Im folgenden Dokument wird ein verbessertes Verfahren zur Bundestagswahl entwickelt. Auch wenn
das vorgestellte Verfahren primär zur Wahl des Deutschen Bundestages gedacht ist, so lässt es
sich doch auf eine Vielzahl anderer Wahlen übertragen. Es wird eigentlich ein allgemeines
Wahlverfahren am Beispiel der Bundestagswahl entwickelt.

Am aktuellen Wahlrecht gibt es verschiedene Kritikpunkte:
\begin{enumerate}
  \item Zwang zu zähneknirschenden Koalitionen mit faulen Kompromissen und vielen Streitigkeiten.
    Da es bei der zersplitterten Parteienlandschaft kaum möglich ist, dass eine Partei die absolute
    Mehrheit gewinnt, müssen Koalitionen gebildet werden. Die Koalitionspartner blockieren sich
    aber oft gegenseitig so sehr, dass nichts herauskommt und alle Partner als Lügner dastehen.

  \item Spitzenkandidaten werden nicht vom Bürger gewählt. Wer eine bestimmte Partei aufgrund
    ihrer Positionen wählt, muss auch ihren Spitzenkandidaten schlucken -- egal wie der ist.

  \item Bei der aktuellen Implementierung des personalisierten Verhältniswahlrechts kann es
    passieren, dass einige Wahlkreise nicht personell im Bundestag repräsentiert sind. Das wurde
    von Friedrich Merz nach der Bundestagswahl 2025 kritisiert. Das vorangegangene System besaß
    dagegen den allgemein kritisierten Nachteil eines immer weiter aufgeblähten Bundestages.
\end{enumerate}

Alle drei Kritikpunkte werden durch das im Folgenden abgeleitete System vermieden.

Der Vorschlag versteht sich als Diskussionsbeitrag in einer noch anzustoßenden
gesamtgesellschaftlichen Diskussion, die mit einer Verfassungsreform oder einer neuen Verfassung
enden könnte.

Unterstützend zur Beschreibung des Verfahrens gibt es eine Simulation als öffentliches
\href{https://github.com/danielschwalbe/ipres}{GitHub-Projekt.} Die enthaltenen Jupyter Notebooks
können dazu verwendet werden, um mit dem Verfahren zu experimentieren und die Wirkung
verschiedener Parameter zu studieren.

Die Simulation ist auch die ultimative Verfahrensreferenz, da sie getestet ist und außerdem
einen Detaillierungsgrad aufweist, den man in einem solchen Konzeptdokument nur sehr schwer
erreichen kann. Hier wird dafür das Warum und Wieso erläutert. Das ist in der Markdown-Dokumentation
zur Simulation nur sehr komprimiert enthalten.
